\documentclass[11pt]{article}
\usepackage[T1]{fontenc}
\usepackage{amsmath,amssymb,geometry}
\geometry{margin=1in}
\begin{document}
\begin{center}
{\Large The Engel correction cannot tend to zero}\\[4pt]
Disproof of Conjecture 00000002760\\[4pt]
ziangni-sys --- October 8, 2026
\end{center}
\paragraph{Scope.}
The problem proposes that the dimension of $k$-Engel linear matrix
spaces is $n^2-n+c(k)$, where $c(k)$ decreases to zero as $k$ grows.
Neither language version excludes $n=1$. We give a counterexample
in this matrix size, which also works if the stated dimension is
interpreted as the maximum possible dimension.

\paragraph{Actual iterated commutators.}
Let $M=M_1(\mathbb Q)$ and take the full linear space $S=M$.
For matrices $X,Y\in M$, define
\[
 E_0(X,Y)=Y,\qquad
 E_{k+1}(X,Y)=[X,E_k(X,Y)]
             =XE_k(X,Y)-E_k(X,Y)X.
\]
These are the usual iterated Engel commutators. Every $1\times1$
matrix is a scalar matrix, so actual matrix multiplication commutes.
It follows immediately that $E_k(X,Y)=0$ for every $k\geq1$ and
all $X,Y\in M$. Thus the full space is $k$-Engel for every positive
$k$. The same conclusion holds if the repeated commutator is
conventionally taken on the other side.

\paragraph{Dimension and maximality.}
The matrix $(1)$ is a basis of $M$ over $\mathbb Q$, so
$\dim_{\mathbb Q}S=1$. Every linear subspace of $M$ has dimension
at most one. Consequently the maximum dimension of a $k$-Engel
linear space in $M_1(\mathbb Q)$ is also exactly one for every
positive $k$.

\paragraph{Contradiction.}
Apply the proposed formula at $n=1$. Since the actual dimension is
one, it forces
\[
 1=1^2-1+c(k)=c(k)\qquad\text{for every }k\geq1.
\]
Hence the correction has a constant tail equal to one and converges
to one. It cannot converge to zero: uniqueness of limits would give
$1=0$. Monotonicity does not remove this contradiction. Therefore
no correction function with the stipulated properties can satisfy
the claimed dimension formula, and the conjecture is false.

\paragraph{Formal certificate.}
The Lean project uses the actual type of $1\times1$ rational matrices,
its genuine multiplication, the recursive iterated commutator, and
the full matrix submodule. It proves the Engel property for every
positive index, calculates its finrank, bounds every subspace's
finrank and certifies the maximal dimension. The final theorem
rules out any real-valued correction satisfying the formula,
monotonicity and convergence to zero. Printed axiom audits use
only the standard axioms of Lean and Mathlib.
\end{document}
